% s-086-p086-m002; original PDF 86, 87
\Aphorism{s-086-p086-m002}{84}
Das Einzige, was ich zu Gunsten des Duells zu sagen wüsste, ist, dass für gewisse verzweifelte Fälle auch ein verzweifeltes Mittel noth thut. Es giebt Sackgassen der Verschuldung und Wirrniss, die keinen Ausweg als den Durchbruch zulassen, Knoten und Unmöglichkeiten, die auf die Gefahr des Blutfliessens hin unter allen Umständen zerschnitten werden müssen. Da bedarf man eines mystischen und seiner Zusammensetzung nach nicht weiter discutiblen Geheimmittels, einer historisch geheiligten Unvernunft, einer groben aber verständlichen Übertragung des inneren Dramas in äussere Drastik. An Stelle der Natur, die hier nicht helfen kann, tritt der feierlich gestabte Unsinn, der supranaturalistisch zu wirken verspricht; wo der Arzt versagt, stellt der Beschwörer, der Teufelsbanner, der Curpfuscher zur rechten Zeit sich ein. Namentlich bei dreieckigen Verhältnissen aller Art (wo, nach Lenau, um ein Mädel zwei Buben sich streiten) wird stets ein Punkt erreicht werden, da die Pistolenmündung zum Mundstück Gottes und berufenen Orakel wird. Aber — und hierin sehe ich eine Widerlegung des gegenwärtigen Duells — ein Arcanum darf nicht in den Apotheken verkauft und von jedem Beliebigen verschrieben werden, das Duell darf nicht gesellschaftliche \Sperr{Forderung} sein! Ein Codex, der den betrogenen Ehemann \Sperr{zwingen} will, mit dem Verführer auf diese Weise abzurechnen, ist eine Barbarei; es muss in jedem Einzelfall dem Betroffenen, unbeschadet seiner Ehre, überlassen bleiben, ob er von dem Geheimmittel Gebrauch machen will oder nicht. Das Duell ist die privateste der Privatsachen, es appellirt an das persönlichste Stück Aberglauben und Rückständigkeit im Menschen: fort mit ihm aus der Öffentlichkeit, fort mit ihm als einem Zubehör der Standesehre, fort mit dem subtilen, steifen, von müssigen Raufbolden erklügelten Talmud von Spitzfindigkeiten, der sich \Foreign{point d’honneur} nennt!
\par

